\subsection*{Logbook Entry 29: Full-universe candidate scan for Hedgehog- and stemness-linked genes}

\textbf{What was being measured.}
A full-universe candidate scan was performed to determine whether Hedgehog-linked, stemness-linked, and SIM2-seed candidate genes were represented anywhere in the selected D3 versus not-D3 methylation probe universe, rather than only in the previously exported directional highlight table.

\textbf{Methods.}
A D3 versus not-D3 probe universe table was identified at
\texttt{/home/satal/cancer\_cells\_hrsm\_sims/results/gse240704/d2\_d3\_probe\_signatures/D3\_vs\_not\_D3\_all\_probes.csv}
and merged with the manifest-derived annotation table
\texttt{results/gse240704/manual\_manifest\_annotation/gpl23976\_annotation\_from\_manifest.parquet}.
Candidate genes included PTCH1, GLI1, GLI2, GLI3, SHH, SOX2, OLIG2, BMP4, PAX6, ASCL1, MYCN, and genes from the SIM2 seed list. Probe-level overlaps were identified by gene-token matching, then summarized by number of detected probes, direction labels when present, and best available absolute effect size.

\textbf{Results.}
Detected candidate genes in the selected universe table were: PAX6, GLI2, BMP4, PTCH1, SHH, OLIG2.
Candidate genes not detected in the same universe table were: ASCL1, GLI1, GLI3, MYCN, SOX2.
This step extends the earlier candidate overlay beyond the exported directional shortlist and therefore provides a stronger basis for interpreting whether the SIM2-centered methylation signal sits inside a broader developmental, Hedgehog-related, or stemness-related methylation pattern.

\textbf{Tables written.}
\texttt{results/gse240704/full\_universe\_candidate\_scan/full\_universe\_candidate\_hits.csv}\\
\texttt{results/gse240704/full\_universe\_candidate\_scan/full\_universe\_candidate\_summary.csv}\\
\texttt{results/gse240704/full\_universe\_candidate\_scan/tables\_tex/table\_full\_universe\_candidate\_scan.tex}\\
\texttt{results/gse240704/full\_universe\_candidate\_scan/full\_universe\_candidate\_summary.json}

\textbf{Figures.}
\texttt{results/gse240704/full\_universe\_candidate\_scan/plots/full\_universe\_candidate\_support.png}

\textbf{Interpretation.}
This step determines whether the earlier absence of Hedgehog- and stemness-linked genes from the exported directional highlight set reflects a true broader absence in the selected D3 versus not-D3 universe, or only a shortlist effect. The interpretation should depend on the detected-versus-not-detected pattern in this broader scan, not on the earlier highlight table alone.

\textbf{Caveats.}
This result still depends on the specific D3 versus not-D3 universe table used for the scan. If the selected table is not the complete probe-universe output, then absence statements remain conditional on that source file.

\textbf{Next steps.}
Use the full-universe scan result to decide whether the manuscript should describe the SIM2-centered signal as largely local and isolated, or whether it sits within a broader candidate-gene methylation program. If the universe table used here is still only partial, the next step is to point this same script to the true complete D3 versus not-D3 differential output.
